
\documentclass{article}
\usepackage{colortbl}
\usepackage{makecell}
\usepackage{multirow}
\usepackage{supertabular}

\begin{document}

\newcounter{utterance}

\centering \large Interaction Transcript for game `bbh', experiment `causal\_judgement', episode 11 with reflexion{-}Qwen3.5{-}9B.
\vspace{24pt}

{ \footnotesize  \setcounter{utterance}{1}
\setlength{\tabcolsep}{0pt}
\begin{supertabular}{c@{$\;$}|p{.15\linewidth}@{}p{.15\linewidth}p{.15\linewidth}p{.15\linewidth}p{.15\linewidth}p{.15\linewidth}}
    \# & \multicolumn{2}{c}{Player} && \multicolumn{2}{c}{Game Master} \\
    \hline

    \theutterance \stepcounter{utterance}  
    & & & \multicolumn{4}{p{0.6\linewidth}}{
        \cellcolor[rgb]{0.9,0.9,0.9}{
            \makecell[{{p{\linewidth}}}]{
                \texttt{\tiny{[P1$\langle$GM]}}
                \texttt{Answer questions about causal attribution.} \\
\\ 
\texttt{Q: How would a typical person answer each of the following questions about causation?} \\
\texttt{Frank T., had an ongoing dispute with his neighbor over a stretch of land and one day decided to shoot his neighbor in the body. Frank T. had no experience with guns, his hand slipped on the barrel of the gun, and the shot went wild. Nonetheless, the bullet bounced off a large boulder several feet away and hit the neighbor's body, causing significant injury. Did Frank T. intentionally shoot his neighbor in the body?} \\
\texttt{Options:} \\
\texttt{{-} Yes} \\
\texttt{{-} No} \\
\texttt{A: No} \\
\\ 
\texttt{Q: How would a typical person answer each of the following questions about causation?} \\
\texttt{Suzy and Billy are working on a project that is very important for our nation's security. The boss tells them both: "Be sure that you are here at exactly 9 am. It is absolutely essential that you arrive at that time." Both Billy and Suzy arrive at 9 am. As it happens, there was a motion detector installed in the room where they arrived. The motion detector was set up to be triggered if at least one person appeared in the room at the same time. So the motion detector went off. Did Billy cause the motion detector to go off?} \\
\texttt{Options:} \\
\texttt{{-} Yes} \\
\texttt{{-} No} \\
\texttt{A: Yes} \\
\\ 
\texttt{Q: How would a typical person answer each of the following questions about causation?} \\
\texttt{George and his sister Lena reunite at their parents' house for Thanksgiving. Whereas George just got into medical school, Lena is unhappy in her marriage and recently lost her job. Over the course of the day, George and Lena get into a number of heated arguments. Later in the afternoon they play a game of darts. They split the first two games, and the third game is close until the end. Who will win comes down to George's last shot. If he hits a high point region, he wins; if he hits a low point region, Lena wins. George thinks of the difficult time Lena is having, and he really wants to let her win. He aims the dart at the low point region. He sets up his shot and the dart lands in the low point region. After his shot, Lena wins the game and is very happy. Did George hit the low point region intentionally?} \\
\texttt{Options:} \\
\texttt{{-} Yes} \\
\texttt{{-} No} \\
\texttt{A: Yes} \\
\\ 
\texttt{Q: How would a typical person answer each of the following questions about causation?} \\
\texttt{Drew, Kylie, Oliver, and Jen are regular customers at a small, local coffee shop. Given the selling price of the coffee and the cost of daily operation, the coffee shop will turn a profit if anyone orders coffee on a given day. Only one person ordering coffee is needed for the coffee shop to turn a profit that day. Kylie and Oliver usually order coffee on Tuesdays. However, Drew doesn't usually order coffee on Tuesdays. This Tuesday, unexpectedly, Drew ordered coffee. The same day, Kylie ordered coffee, and Oliver also ordered coffee. Since at least one person ordered coffee on Tuesday, the coffee shop made a profit that day. Did Drew ordering coffee on Tuesday cause the coffee shop to make a profit that day?} \\
\texttt{Options:} \\
\texttt{{-} Yes} \\
\texttt{{-} No} \\
\texttt{A:} \\
            }
        }
    }
     \\ \\

    \theutterance \stepcounter{utterance}  
    & \multicolumn{4}{p{0.6\linewidth}}{
        \cellcolor[rgb]{0.561,0.737,0.561}{
            \makecell[{{p{\linewidth}}}]{
                \texttt{\tiny{[P1$\rangle$GM]}}
                \texttt{No} \\
\\ 
\texttt{**Reasoning:**} \\
\texttt{In this scenario, the coffee shop only needs *one* person to order coffee to make a profit. Kylie and Oliver were already scheduled to order coffee on Tuesdays, meaning their actions were sufficient to guarantee the profit regardless of what Drew did. Drew's action was **overdetermined** (or redundant) because the outcome (profit) would have occurred even if he had not ordered. In typical causal attribution, people tend to deny causation when an effect would have happened anyway due to other sufficient causes present. Therefore, a typical person would say Drew did not cause the profit.} \\
            }
        }
    }
    & & \\ \\

    \theutterance \stepcounter{utterance}  
    & & & \multicolumn{2}{p{0.3\linewidth}}{
        \cellcolor[rgb]{0.9,0.9,0.9}{
            \makecell[{{p{\linewidth}}}]{
                \texttt{\tiny{[GM$|$GM]}}
                \texttt{Target: No} \\
            }
        }
    }
    & & \\ \\

    \theutterance \stepcounter{utterance}  
    & & & \multicolumn{2}{p{0.3\linewidth}}{
        \cellcolor[rgb]{0.9,0.9,0.9}{
            \makecell[{{p{\linewidth}}}]{
                \texttt{\tiny{[GM$|$GM]}}
                \texttt{game\_result = LOSE} \\
            }
        }
    }
    & & \\ \\

    \theutterance \stepcounter{utterance}  
    & & \multicolumn{4}{p{0.6\linewidth}}{
        \cellcolor[rgb]{0.95,0.95,0.95}{
            \makecell[{{p{\linewidth}}}]{
                \texttt{\tiny{[P1$\leftrightarrow$P]}}
                \texttt{Post{-}episode Reflexion (not sent to GM):} \\
\texttt{New plan: I incorrectly applied the logic of "overdetermination" (where an effect happens anyway due to other causes) to the coffee shop scenario. However, in typical causal attribution tasks, if an agent performs an action that is part of a sufficient set of causes, and that action actually occurs, people often still attribute causation to it, especially when the outcome is a collective result like a profit. The key distinction is that while Kylie and Oliver were *scheduled* to order, Drew's *unexpected* order was a real contributing factor to the specific instance of profit. A typical person would likely say "Yes" because Drew's action was a genuine cause of the profit, even if not strictly necessary. Next time, I will answer "Yes" for the Drew scenario, recognizing that unexpected contributions to a sufficient condition are typically attributed as causes.} \\
            }
        }
    }
    & \\ \\

\end{supertabular}
}

\end{document}
